
\documentclass{article}
\usepackage{colortbl}
\usepackage{makecell}
\usepackage{multirow}
\usepackage{supertabular}

\begin{document}

\newcounter{utterance}

\centering \large Interaction Transcript for game `adventuregame', experiment `potion\_brewing\_basic\_undefined', episode 1 with reka{-}flash{-}3.1{-}t0.0.
\vspace{24pt}

{ \footnotesize  \setcounter{utterance}{1}
\setlength{\tabcolsep}{0pt}
\begin{supertabular}{c@{$\;$}|p{.15\linewidth}@{}p{.15\linewidth}p{.15\linewidth}p{.15\linewidth}p{.15\linewidth}p{.15\linewidth}}
    \# & \multicolumn{2}{c}{Player} && \multicolumn{2}{c}{Game Master} \\
    \hline

    \theutterance \stepcounter{utterance}  
    & & & \multicolumn{4}{p{0.6\linewidth}}{
        \cellcolor[rgb]{0.9,0.9,0.9}{
            \makecell[{{p{\linewidth}}}]{
                \texttt{\tiny{[P1$\langle$GM]}}
                \texttt{You are playing a text adventure game. I will describe what you can perceive in the game. You write the single action you want to take in the game starting with >. Only reply with an action.} \\
\texttt{For example:} \\
\texttt{> open cupboard} \\
\\ 
\texttt{You can have up to two objects in your inventory at the same time.} \\
\\ 
\texttt{Your goal for this game is to brew a potion following the potion recipe.} \\
\\ 
\texttt{Once you have achieved your goal, write "> done" to end the game.} \\
\\ 
\texttt{You are in a bedroom now. There are a wardrobe, a potion recipe, a bed, a book and a pillow. The wardrobe is closed. The potion recipe is on the bed. The book is on the bed. The pillow is on the floor. There is a passage to a reading room here.} \\
            }
        }
    }
     \\ \\

    \theutterance \stepcounter{utterance}  
    & \multicolumn{4}{p{0.6\linewidth}}{
        \cellcolor[rgb]{0.561,0.737,0.561}{
            \makecell[{{p{\linewidth}}}]{
                \texttt{\tiny{[P1$\rangle$GM]}}
                \texttt{<reasoning>} \\
\texttt{I need to brew a potion following the recipe. Let me check what's in the room. The description says there's a wardrobe, potion recipe on the bed, a book, pillow, and a passage to the reading room. The wardrobe is closed. The inventory can hold up to two objects.} \\
\\ 
\texttt{First, I should probably look at the potion recipe to know what ingredients I need. The recipe is on the bed, so maybe I need to pick it up. But wait, the bed has both the recipe and a book. If I take the recipe, maybe I can read it. Let me check the possible actions. The user can take objects, so maybe "take potion recipe" or "take book". But the goal is to brew the potion, so the recipe is essential. Let me start by taking the recipe. That way I can see what's needed.} \\
\\ 
\texttt{So action: > take potion recipe. Then I can read it. But wait, the bed has both the recipe and the book. If I take the recipe, the book stays. But maybe the book is a distractor. Alternatively, maybe the book is needed for something else. But the goal is the potion, so focus on the recipe. Once I have the recipe, I need to gather ingredients. The room has a wardrobe. Maybe the wardrobe contains ingredients. Let me check. The wardrobe is closed, so I need to open it first. So steps: take recipe, open wardrobe, take ingredients from wardrobe, then go to the reading room? Wait, the passage is to the reading room. Maybe the reading room has a cauldron or something needed for brewing. But the current room has a bed, wardrobe, etc. Let me think.} \\
\\ 
\texttt{Wait, the initial description says "You are in a bedroom now. There are a wardrobe, a potion recipe, a bed, a book and a pillow. The wardrobe is closed. The potion recipe is on the bed. The book is on the bed. The pillow is on the floor. There is a passage to a reading room here." So the bedroom has the wardrobe, bed, book, pillow, and a passage. The recipe is on the bed. So first, take the recipe. Then, check the wardrobe. Open the wardrobe to see if there are ingredients. Let's assume the wardrobe has ingredients like herbs, vials, etc. Once I have the recipe, I can see what's needed. For example, maybe} \\
            }
        }
    }
    & & \\ \\

    \theutterance \stepcounter{utterance}  
    & & & \multicolumn{2}{p{0.3\linewidth}}{
        \cellcolor[rgb]{0.9,0.9,0.9}{
            \makecell[{{p{\linewidth}}}]{
                \texttt{\tiny{[GM$|$GM]}}
                \texttt{Round: 0} \\
            }
        }
    }
    & & \\ \\

    \theutterance \stepcounter{utterance}  
    & & & \multicolumn{2}{p{0.3\linewidth}}{
        \cellcolor[rgb]{0.9,0.9,0.9}{
            \makecell[{{p{\linewidth}}}]{
                \texttt{\tiny{[GM$|$GM]}}
                \texttt{command\_tag\_missing} \\
            }
        }
    }
    & & \\ \\

    \theutterance \stepcounter{utterance}  
    & & & \multicolumn{2}{p{0.3\linewidth}}{
        \cellcolor[rgb]{0.9,0.9,0.9}{
            \makecell[{{p{\linewidth}}}]{
                \texttt{\tiny{[GM$|$GM]}}
                {'goal_states_achieved': [], 'game_successfully_finished': False}
            }
        }
    }
    & & \\ \\

    \theutterance \stepcounter{utterance}  
    & & & \multicolumn{2}{p{0.3\linewidth}}{
        \cellcolor[rgb]{0.9,0.9,0.9}{
            \makecell[{{p{\linewidth}}}]{
                \texttt{\tiny{[GM$|$GM]}}
                \texttt{command\_tag\_missing} \\
            }
        }
    }
    & & \\ \\

\end{supertabular}
}

\end{document}
